\chapter{제안 방법}\label{ch:method}

\section{문제 정의}
\begin{definition}[문제]\label{def:problem}
  입력 $\mathbf{x} \in \mathbb{R}^d$ 에 대해 출력 $y$ 를 예측하는 함수 $f$ 를 찾는다.
\end{definition}

\section{제안 모델}
목적 함수는 \cref{eq:objective}와 같습니다.
\begin{equation}\label{eq:objective}
  \min_{\theta} \; \frac{1}{N}\sum_{i=1}^{N} \mathcal{L}\bigl(f_\theta(\mathbf{x}_i), y_i\bigr) + \lambda \lVert \theta \rVert_2^2
\end{equation}

\begin{theorem}\label{thm:main}
  적절한 조건 하에서 \cref{eq:objective}의 해는 유일하다.
\end{theorem}
\begin{proof}
  증명을 서술합니다.
\end{proof}

\section{델파이 조사 설계 및 분석 방법}\label{sec:delphi}

VR 검도훈련 콘텐츠 구성요소의 타당성은 델파이 기법~\cite{dalkey1963delphi}으로 검증하였다.
전문가 패널은 검도 지도자, 체육교육·스포츠과학 연구자, VR·게임 개발자, 교육공학 전문가로 구성하고,
동일한 패널을 대상으로 조사를 반복하며 이전 차수의 중앙값과 사분위 범위를 제시하여 의견 수렴을 유도하였다.
설문 문항은 \cref{app:survey}와 같다.

각 문항은 내용타당도 비율(CVR)~\cite{lawshe1975cvr}, 합의도, 수렴도~\cite{lee2001delphi}, 변이계수(CV)로 분석하였다.
응답자 수를 $N$, 4점 이상으로 응답한 전문가 수를 $n_e$, 중앙값을 $Mdn$, 제1·3사분위수를 $Q_1$, $Q_3$라 할 때
\begin{align}
  \mathrm{CVR} &= \frac{n_e - N/2}{N/2}, &
  \text{합의도} &= 1 - \frac{Q_3 - Q_1}{Mdn}, &
  \text{수렴도} &= \frac{Q_3 - Q_1}{2}, &
  \mathrm{CV} &= \frac{SD}{M}
  \label{eq:delphi}
\end{align}
이다. \Cref{tab:delphi-criteria}의 기준을 모두 충족한 문항을 채택하고,
하나라도 충족하지 못한 문항은 전문가 의견을 반영하여 수정·삭제 여부를 검토하였다.

\begin{table}[htbp]
  \centering
  \caption{델파이 조사 문항 채택 기준}\label{tab:delphi-criteria}
  \begin{tabular}{lll}
    \toprule
    지표 & 기준 & 의미 \\
    \midrule
    CVR & $\geq$ 최소값 ($N=12$일 때 .56) & 내용타당도 \\
    합의도 & $\geq .75$ & 전문가 간 의견 일치 \\
    수렴도 & $\leq .50$ & 응답의 집중 정도 \\
    CV & $\leq .50$ & 응답의 안정성 \\
    \bottomrule
  \end{tabular}
\end{table}

\section{알고리즘}
\begin{algorithm}[htbp]
  \caption{제안 알고리즘}\label{alg:proposed}
  \KwIn{데이터 $\{(\mathbf{x}_i, y_i)\}_{i=1}^N$, 학습률 $\eta$}
  \KwOut{학습된 파라미터 $\theta$}
  $\theta \leftarrow \theta_0$\;
  \For{$t = 1$ \KwTo $T$}{
    $\theta \leftarrow \theta - \eta \nabla_\theta \mathcal{L}$\;
  }
  \Return{$\theta$}\;
\end{algorithm}
